\noindent\textbf{Scope and credit.}
The first two sections develop bounded parts of the class C problems
\hyperref[prob:GP-1986-33]{GP-1986-33} and
\hyperref[prob:GP-1986-60]{GP-1986-60}. The last two are verification
annotations to questions already answered in the literature or notebook.
The markers \ours\ and \oursComputed\ identify proofs and computations
supplied here; they do not assert historical novelty. No finite calculation
below is promoted to an all-level spectral theorem or a large-\(N\) result.

\subsection{Quantum KdV: a level-two calculation and a Jordan obstruction}
\label{dev:kdv}

\paragraph{The question and its scope.}
At \GSpage{P001318} the notebook asks to quantize the KdV integrals
and then compute the spectrum in irreducible Virasoro representations.
Reading the latter as the spectrum of the quantized integrals, rather
than of \(L_0\) alone, is an editorial reconstruction based on the
preceding clause. The notes specify no central charge, highest weight,
positivity, or completeness criterion. A precise modern target fixes
\(c,h\in\mathbb C\), the irreducible highest-weight module \(L(c,h)\),
charge normalizations, and a grade \(n\); one seeks the joint generalized
eigenspaces and nilpotent actions, including multiplicities. An ordinary
joint eigenbasis requires an additional semisimplicity assertion.

The construction of commuting charges and their spectral classification
are different problems.
\citet[eq.~(11)]{bazhanovlukyanovzamolodchikov1996integrable}
gives the normalization used below (December 1994 preprint; 1996
publication). The opened May 2026 version of
\citet[Definitions~6.5--6.6 and Theorem~6.8]{masoeroruzza2026odeim}
proves a completeness statement at \(c=-2\),
\(h=(\beta^2-1)/8\), \(\beta\notin\mathbb Z\), for the specified
admissible QQ solutions modulo their entire-function gauge.
This is not a theorem for every \(c,h\). The calculation below
does not depend on that preprint.

\paragraph{Conventions.}
Let \(M(c,h)\) be the Verma module generated by \(v\), where
\[
 L_mv=0\ (m>0),\quad L_0v=hv,\quad
 [L_m,L_n]=(m-n)L_{m+n}
       +\frac c{12}(m^3-m)\delta_{m+n,0}.
\]
Its grade-two PBW basis is \(v_1=L_{-2}v,\ v_2=L_{-1}^2v\). Set
\[
 I_3=2\sum_{n\geq1}L_{-n}L_n+
       L_0^2-\frac{c+2}{12}L_0+\frac{c(5c+22)}{2880}.
\]
The sum is finite on every vector of the algebraic module.

\begin{proposition}[The level-two charge]\ours\ 
For
\[
 q_2=(h+2)^2-\frac{(c+2)(h+2)}{12}+\frac{c(5c+22)}{2880},
\]
the matrix of \(I_3\) on \(M(c,h)[2]\) is
\[
 q_2\,1+\begin{pmatrix}8h+c&12h\\6&8h+4\end{pmatrix}.
\]
Its characteristic roots, with algebraic multiplicity, are
\[
 q_2+8h+\frac{c+4}{2}
       \ \pm\ \frac12\sqrt{(c-4)^2+288h}.
\]
It is diagonalizable on this two-dimensional subspace exactly when
\((c-4)^2+288h\ne0\).
\end{proposition}
\begin{proof}
The commutators give
\[
\begin{aligned}
 L_1v_1&=3L_{-1}v,&L_1v_2&=(4h+2)L_{-1}v,\\
 L_2v_1&=(4h+c/2)v,&L_2v_2&=6hv.
\end{aligned}
\]
For example,
\(L_1L_{-1}^2v=2L_0L_{-1}v+2L_{-1}L_0v
=(4h+2)L_{-1}v\).
All \(L_n\), \(n>2\), annihilate this grade. Multiplication by
\(2L_{-1}\) and \(2L_{-2}\) gives the matrix.
Its nonscalar part \(B\) has trace \(16h+c+4\) and discriminant
\((\operatorname{tr}B)^2-4\det B=(c-4)^2+288h\).
If this vanishes, the repeated-root matrix is not scalar, since its
lower-left entry is \(6\); hence it is not diagonalizable.
\end{proof}

\begin{proposition}[A Jordan block in an irreducible quotient]\ours\ 
For \(c=1,\ h=-1/32\), grade two survives unchanged in \(L(c,h)\).
On that grade,
\[
 I_3=\frac{28893}{5120}\,1+
       \begin{pmatrix}-3/2&-3/8\\6&3/2\end{pmatrix}.
\]
The second matrix \(N\) is nonzero and satisfies \(N^2=0\).
Thus there is no ordinary simultaneous eigenbasis for the local charges
on this irreducible module.
\end{proposition}
\begin{proof}
Substitution gives the displayed matrix and \(N^2=0\).
Every submodule is graded: it is invariant under the diagonal operator
\(L_0\), and each vector has only finitely many weight components.
A proper submodule meeting grade two would therefore have a nonzero
singular vector in grade one or two. In grade one the condition is
\(2h=0\), which fails. In grade two, annihilation by \(L_1,L_2\)
requires
\[
 \begin{pmatrix}3&4h+2\\4h+c/2&6h\end{pmatrix}\binom ab=0.
\]
At the stated parameters its determinant is \(-81/64\ne0\).
There is consequently no such singular vector, and no proper
submodule meets grade two. The Jordan block persists in the
irreducible quotient. A simultaneous eigenbasis would diagonalize
\(I_3\), a contradiction.
\end{proof}

\paragraph{Interpretation and checks.}
This example is nonunitary: the highest-weight contravariant form
is negative on \(L_{-1}v\). It does not challenge finite-grade
diagonalizability with a positive-definite unitary inner product.
It also does not refute a notebook conjecture: the notebook asks for
a spectrum, without asserting semisimplicity.

The independent program \texttt{code/root/check\_\allowbreak{}quantum\_\allowbreak{}kdv.py}
constructs the action by recursive PBW rewriting, without entering
the proposed matrix. A separate Gram calculation gives
\[
 G=\begin{pmatrix}4h+c/2&6h\\6h&4h(2h+1)\end{pmatrix},
 \quad
 \det G=2h\{16h^2+(2c-10)h+c\}.
\]
It checks \(GI_3=I_3^{\mathsf T}G\); at the exceptional parameters
\(\det G=-81/1024\). The program also checks Gram symmetry in
grades one, three, and four at two other rational parameter choices.
All exact checks passed; the output is supplied.

\paragraph{Remaining obstruction.}
This calculation does not organize arbitrary grades, settle every
singular-vector quotient, or prove ODE/IM completeness. The reviewed
sources establish particular constructions and regimes, without
documenting a reason that the general problem was abandoned.
A further bounded run could construct the charge algebra in grades
three through six and separate its Jordan loci from degeneration
of the Verma form. It would give a finite classification, not an
all-grade theorem.

\subsection{The multi-source unitary integral: evaluation, error control, and a failed extension}
\label{dev:unitary}

\paragraph{Supported and reconstructed formulations.}
On \GSpage{P001330}, Kazakov asks for the integral with two diagonal
matrix pairs, writing identical pairs first. We use probability
Haar measure, \(\int_{U(N)}dU=1\), and write
\[
 Z_N(P,Q)=\int_{U(N)}
 \exp\!\left(\sum_{\mu=1}^r
   \operatorname{Tr}(U^\ast P^\mu UQ^\mu)\right)dU,\qquad
 P^\mu=\operatorname{diag}(p_i^\mu),\quad
 Q^\mu=\operatorname{diag}(q_j^\mu).
\]
The notebook supports \(r=2\); arbitrary finite \(r\), complex entries,
and the explicit Haar normalization are editorial extensions.
Put \(a_{ij}=\sum_\mu p_i^\mu q_j^\mu\). The exponent becomes
\(X_A(U)=\sum_{ij}a_{ij}|U_{ij}|^2\).
The diagonal entries are fixed parameters, so no contour for them is
needed.

\paragraph{Literature and the remaining target.}
The \(N=2\) reduction and the failure of an immediate localization
extension to \(N>2\) are treated in
\citet[Sections~2.2--2.3 and Appendix~A]{holguin2024multimatrix}
(July 2023 preprint; 2024 publication). Their displayed angular measure
and our probability Haar measure must be compared after fixing the
overall normalization. All-rank polynomial Haar moments are given by
\citet[Corollary~2.4]{collinssniady2006integration}
(February 2004 preprint; 2006 publication).
Thus an assertion that the integral has never been evaluated would
be wrong if evaluation includes a convergent explicit series.
A meaningful remaining target is a tractable formula or asymptotic
description for general multi-source data in a stated regime.

\begin{proposition}[The normalized \(U(2)\) integral]\ours\ 
For \(A\in M_2(\mathbb C)\), put
\(\alpha=a_{11}+a_{22}\), \(\beta=a_{12}+a_{21}\).
Then
\[
 Z_2(A)=
 \begin{cases}
 (e^\alpha-e^\beta)/(\alpha-\beta),&\alpha\ne\beta,\\
 e^\alpha,&\alpha=\beta.
 \end{cases}
\]
These are the values of one entire function. Equivalently,
\(Z_2=e^s\operatorname{sinhc}(z)\), where
\(s=(\alpha+\beta)/2\), \(z=(\alpha-\beta)/2\),
\(\operatorname{sinhc}(z)=\sinh(z)/z\) for \(z\ne0\), and
\(\operatorname{sinhc}(0)=1\).
\end{proposition}
\begin{proof}
For Haar \(U\in U(2)\), \(t=|U_{11}|^2\) is uniform on \([0,1]\).
Indeed a first column is obtained by normalizing two independent
complex Gaussians. Their squared moduli are independent exponentials.
The change \((x,y)=(rt,r(1-t))\), with Jacobian \(r\), gives
density \(\int_0^\infty re^{-r}\,dr=1\) for \(t\).
Unitarity then yields
\[
 (|U_{ij}|^2)=\begin{pmatrix}t&1-t\\1-t&t\end{pmatrix},
 \qquad Z_2=\int_0^1e^{\beta+(\alpha-\beta)t}\,dt.
\]
Integration gives the formula; uniform convergence of its exponential
series on compact parameter sets gives entireness and the diagonal value.
\end{proof}

\begin{proposition}[All-rank convergent evaluation with an explicit tail]\ours\ 
For \(A\in M_N(\mathbb C)\), set
\(M=\sum_i\max_j|a_{ij}|\). Then
\[
 Z_N(A)=\sum_{m=0}^\infty\frac{\mathbb E X_A^m}{m!},\qquad
 \left|Z_N(A)-\sum_{m=0}^K\frac{\mathbb E X_A^m}{m!}\right|
 \leq\frac{e^M M^{K+1}}{(K+1)!}.
\]
Each moment is an explicit finite sum, even for \(m>N\).
Index a matrix by \(S_m\), and put
\[
 G_{\sigma,\tau}=N^{\#\operatorname{cycles}(\sigma^{-1}\tau)},\qquad
 d_{\boldsymbol i}(\sigma)=
     \prod_{k=1}^m\delta_{i_k,i_{\sigma(k)}}.
\]
Writing \(G^+\) for the Moore--Penrose inverse, for \(m\geq1\) one has
\[
 \mathbb E X_A^m=
 \sum_{\boldsymbol i,\boldsymbol j\in\{1,\ldots,N\}^m}
 \left(\prod_{k=1}^m a_{i_kj_k}\right)
 \sum_{\sigma,\tau\in S_m}
 d_{\boldsymbol i}(\sigma)(G^+)_{\sigma,\tau}
 d_{\boldsymbol j}(\tau).
\]
For rational \(A\), all moments are rational.
\end{proposition}
\begin{proof}
Each row of \((|U_{ij}|^2)\) sums to one, giving \(|X_A|\leq M\).
The exponential series converges absolutely and uniformly on the
compact integration group. Termwise integration and the scalar bound
\[
 \left|e^z-\sum_{m=0}^Kz^m/m!\right|
 \leq e^{|z|}|z|^{K+1}/(K+1)!
\]
prove the series and error bound.

On \(\mathcal H=(\mathbb C^N)^{\otimes m}\), averaging conjugation
by \(U^{\otimes m}\) is the Hilbert--Schmidt orthogonal projection
onto the commutant. Schur--Weyl duality says that this commutant is
spanned by the permutation operators \(P_\sigma\).
Their Gram matrix is
\(\operatorname{Tr}(P_\sigma^\ast P_\tau)=G_{\sigma,\tau}\).
The orthogonal projection onto a possibly dependent family uses its
Gram pseudoinverse. Its diagonal coefficients give
\[
 \mathbb E\prod_{k=1}^m|U_{i_kj_k}|^2
 =\sum_{\sigma,\tau}
 d_{\boldsymbol i}(\sigma)(G^+)_{\sigma,\tau}
 d_{\boldsymbol j}(\tau).
\]
Expansion of \(X_A^m\) proves the stated formula.
The pseudoinverse of the integer matrix \(G\) is rational, by
rational rank factorization and inversion of nonsingular rational
Gram matrices. This proves the rationality assertion. A usual
inverse would be invalid when the permutation operators are
dependent.
\end{proof}

\paragraph{Independent low-degree identities.}\ours\ 
Let \(S=\sum_{ij}a_{ij}\), \(F=\sum_{ij}a_{ij}^2\),
\(R=\sum_i(\sum_j a_{ij})^2\), and
\(C=\sum_j(\sum_i a_{ij})^2\). For \(N\geq2\),
\[
 \mathbb EX_A=S/N,\qquad
 \mathbb EX_A^2=
 \frac{S^2+F}{N^2-1}-\frac{R+C}{N(N^2-1)}.
\]
The squares are algebraic, not absolute squares.
For \(N=2\), independent integration gives
\[
 \mathbb EX_A^m=\frac1{m+1}\sum_{j=0}^m\alpha^j\beta^{m-j},
\]
including coincident endpoints.

\begin{example}[Failure of a specific permutation-only extension]\oursComputed\ 
The following stronger formula is constructed for this audit.
The notebook does not assert that its saddle approximation is exact.
Take the two-component diagonal data to be the rows of
\[
 P=\begin{pmatrix}0&0\\1&0\\0&1\end{pmatrix},\quad
 Q=\begin{pmatrix}0&0\\2&1\\1&3\end{pmatrix},\quad
 A=PQ^{\mathsf T}=
   \begin{pmatrix}0&0&0\\0&2&1\\0&1&3\end{pmatrix}.
\]
For \(\sigma\in S_3\), define
\[
 S_\sigma=\sum_i a_{i,\sigma(i)},\quad
 D_\sigma=\prod_{i<j}
  (a_{i,\sigma(i)}+a_{j,\sigma(j)}
                  -a_{i,\sigma(j)}-a_{j,\sigma(i)}).
\]
The naive extension of the ordinary rank-one fixed-point expression is
\[
 F(t)=\frac2{t^3}\sum_{\sigma\in S_3}
                   \frac{e^{tS_\sigma}}{D_\sigma}.
\]
Its exact data are
\[
\begin{array}{c|rrrrrr}
 \sigma&123&132&213&231&312&321\\ \hline
 S_\sigma&5&2&3&1&1&2\\
 D_\sigma&18&-3&-8&-2&-2&-3
\end{array}
\]
so that
\[
 \sum_\sigma D_\sigma^{-1}=-125/72,\qquad
 F(t)=-\frac{125}{36}t^{-3}+O(t^{-2}).
\]
The true probability-normalized Haar integral is entire and has
\[
 Z_3(tA)=1+\frac73t+\frac{71}{24}t^2+\frac{65}{24}t^3+O(t^4).
\]
Consequently this permutation-only expression cannot be the general
two-source formula.
\end{example}
\begin{proof}
The six substitutions give the table. All denominator factors are
nonzero. Their interpretation is checked independently: rotating
rows \(i,j\) of the permutation matrix by a real Givens rotation
changes the phase to
\[
 S_\sigma-
 (a_{i,\sigma(i)}+a_{j,\sigma(j)}
      -a_{i,\sigma(j)}-a_{j,\sigma(i)})\sin^2\theta.
\]
The code differentiates all eighteen curves to verify the quadratic
factors. Summing the reciprocals gives \(-125/72\), proving the pole.
The genuine integral is entire by the preceding proposition, and
its first moments are computed independently from the Gram formula.
\end{proof}

\paragraph{Verification and practical limits.}
The program \texttt{check\_\allowbreak{}unitary\_\allowbreak{}integral.py} compares the exact
Gram calculation with the \(U(2)\) integral for real, complex, and
coincident endpoints and with the ordinary rank-one HCIZ determinant
at \(N=3\). It tests \(m>N\): the Gram ranks are \(5\) for
\((N,m)=(2,3)\) and \(14\) for \((2,4)\).
All comparisons pass.
For \(A=\left(\begin{smallmatrix}1&2\\3&5\end{smallmatrix}\right)\),
centering at \(11/2\) leaves radius \(1/2\).
At order 20 the computed error is about \(2.26\cdot10^{-27}\),
below the analytic bound \(3.77\cdot10^{-24}\).
These decimals illustrate the proved bound; they are not an
interval-arithmetic certificate.
The separate \texttt{check\_\allowbreak{}stationary\_\allowbreak{}formula.py} checks the
counterexample by both the quadratic variations and exact moments.

The procedure is convergent but factorial in degree. It does not
supply an efficient large-\(N\) algorithm or a compact generic
formula. A further bounded run could exploit permutation symmetries,
compute a prescribed set of \(N=3\) moments, and test a precisely
stated approximation with rigorous tails. The later literature
explicitly revisits the problem, so there is no basis to call it
abandoned.

\subsection{An annotation on simultaneous complex structures}
\label{dev:complex}

\paragraph{Source and later result.}
Gelfand's question at \GSpage{P001276} asks for a complex structure
making all the real codimension-two subspaces complex. The notebook
already warns that even-dimensional intersections are insufficient.
Here a complex structure means a single \emph{linear}
\(J\in\operatorname{End}_{\mathbb R}(\mathbb R^4)\) with \(J^2=-1\).
Nonlinear equivalence of complements is a different question.
\citet[Theorems~1 and~3]{batzies2007planes} gives the four-plane
classification and the transverse normal form (received 2001,
revised 2005, published 2007). The following is an elementary
rederivation, not a new classification.

\begin{proposition}[A finite-family criterion in the transverse case]\ours\ 
Let \(P_1,\ldots,P_k\), \(k\geq3\), be real two-dimensional linear
subspaces of \(\mathbb R^4\), with \(P_i\cap P_j=\{0\}\) for
\(i\ne j\). There are real coordinates with
\[
 P_1=\mathbb R^2\oplus0,\quad P_2=0\oplus\mathbb R^2,\quad
 P_3=\operatorname{graph}(1),\quad
 P_j=\operatorname{graph}(A_j)\ (j\geq4).
\]
A common linear complex structure exists exactly when either every
\(A_j\) is scalar, or a nonscalar \(A\) among them satisfies
\[
 (\operatorname{tr}A)^2-4\det A<0,\qquad
 [A,A_j]=0\quad\hbox{for every }j.
\]
For \(k=4\), the condition is that \(A_4\) is scalar or has negative
discriminant.
\end{proposition}
\begin{proof}
In \(P_1\oplus P_2\), transversality makes \(P_3\) the graph of an
isomorphism \(T:P_1\to P_2\). A coordinate change in \(P_2\)
makes \(T=1\); every other plane is then a graph.
The hypotheses imply
\(\det A_j\ne0\), \(\det(A_j-1)\ne0\), and
\(\det(A_i-A_j)\ne0\) for distinct graph planes.

A linear \(J\) preserving \(P_1,P_2\) is block diagonal.
Preserving \(\operatorname{graph}(1)\) makes its blocks equal:
\(J=\operatorname{diag}(B,B)\). The remaining conditions are
\(B^2=-1\) and \(BA_j=A_jB\).
In a basis \(e,Be\), the centralizer of \(B\) consists exactly of
\(a\,1+bB\). A nonscalar such matrix has discriminant \(-4b^2<0\),
and all such matrices commute. This proves necessity.
Conversely, if every \(A_j\) is scalar, any complex structure \(B\)
works. Otherwise, for a chosen \(A\) with discriminant \(\Delta<0\),
Cayley--Hamilton gives
\[
 B=\frac{2A-(\operatorname{tr}A)1}{\sqrt{-\Delta}},
 \qquad B^2=-1.
\]
Every matrix commuting with \(A\) commutes with \(B\).
Thus \(\operatorname{diag}(B,B)\) has all the required invariances.
\end{proof}

\begin{corollary}[The intersection lattice is insufficient]\ours\ 
Arrangements of four pairwise transverse two-planes in \(\mathbb R^4\)
can have the same intersection lattice, and be joined by a transverse
continuous deformation, while only one admits a common linear
complex structure.
\end{corollary}
\begin{proof}
Keep the first three planes in the normal form and use
\(A_t=\operatorname{diag}(2,2+t)\), \(0\leq t\leq1\).
The determinants \(2(2+t)\) and \(1+t\) of \(A_t\) and \(A_t-1\)
stay positive. All pairwise intersections stay zero and the
intersection lattice stays the rank-two uniform lattice.
At \(t=0\) the matrix is scalar; at \(t>0\) its discriminant is
\(t^2>0\). Apply the proposition.
\end{proof}

\paragraph{Independent check and boundary of the result.}
The code \texttt{check\_\allowbreak{}complex\_\allowbreak{}structures.py} independently solves
the full four-by-four invariance equations. It checks scalar,
nonreal, distinct-real, and Jordan cases and commuting and
noncommuting families. In the last case, the invariant algebra is
only the real scalars, excluding \(J^2=-1\) directly. All checks
agree. This settles the stated transverse four-dimensional linear
subcase; it is not a combinatorial characterization of arbitrary
higher-dimensional or nonlinear configurations. No unspecified
extension is labelled an open conjecture here.

\subsection{An exact audit of the cube's 96 regions}
\label{dev:cube}

\paragraph{The recorded answer.}
The Gelfand--T.~Alekseevskaya question at \GSpage{P001364},
item \hyperref[item:1986-706.3]{1986-706.3}, already has the written
answer \(96\). We count the connected components of the
\emph{interior} of a cube after removing every distinct affine
plane determined by three vertices. This is verification of the
recorded answer, not a new problem.

\begin{proposition}[The cube region count]\oursComputed\ 
For \([-1,1]^3\), the 56 vertex triples determine 20 distinct planes.
Six are faces and do not cut the interior. The other 14 planes cut
the interior into exactly 96 regions.
\end{proposition}
\begin{proof}
The cross-product calculation in
\texttt{check\_\allowbreak{}cube\_\allowbreak{}regions.py} enumerates and deduplicates the planes.
The interior-cutting ones are
\[
 x=\pm y,\quad x=\pm z,\quad y=\pm z,\qquad
 \varepsilon_1x+\varepsilon_2y+\varepsilon_3z=1
     \quad(\varepsilon_i\in\{-1,1\}).
\]
The last eight equations describe distinct planes.
The code computes every triple intersection of the full 20-plane
family in the closed cube and includes its vertices, giving
35 candidate points. Every vertex of every intermediate cell
belongs to this set: a vertex of a full-dimensional bounded
three-dimensional polyhedron lies on at least three independent
bounding planes.

Process the 14 cutting planes. A convex cell splits exactly when
its candidate points have both a strict positive and a strict
negative value of the new plane equation. Linear functions attain
their extrema at vertices, so this is necessary and sufficient.
Points on the plane are kept in both new closed cells.
Exact rational arithmetic gives the successive counts
\[
 1,2,4,6,8,12,15,18,26,36,46,56,72,84,96.
\]
The supplied \texttt{cube\_\allowbreak{}witnesses.json} gives a rational interior
point of each final cell. All 96 sign vectors are distinct, with
no zero coordinate. The cells are convex and hence connected.
\end{proof}

\paragraph{Independent count.}
The second calculation uses the intersection poset rather than
cell splitting. For an arrangement restricted to an open convex
set, the region count is the sum of \(|\mu|\) over flats meeting
that set. The usual formula is stated in
\citet[Theorem~A, p.~100]{greenezaslavsky1983whitney}.
The restriction follows by inclusion--exclusion with compactly
supported Euler characteristic: each nonempty flat--cube
intersection is open convex in its affine span, contributing
\((-1)^{\dim}\); each region contributes \((-1)^3\).
Grouping identical intersections gives the Möbius sum, with its
alternating codimension signs.

Exact interval tests find 43 line flats meeting the cube interior,
with total \(|\mu|=47\), and 21 point flats with total \(|\mu|=34\).
Thus the second count is
\[
                    1+14+47+34=96.
\]
Flats only meeting the boundary are excluded. Both independent
counts agree with the notebook; no correction to its answer is
needed.



